\clearpage
\chapter*{Đính chính cho Chuyên đề số 221}
\addcontentsline{toc}{chapter}{Đính chính cho Chuyên đề số 221}
\setcounter{footnote}{3}
\sid{e221}
\markboth{Đính chính 221}{Đính chính 221}
\src{302}

\noindent\textbf{6. TDTE IV: Các lược đồ Hilbert.}

\begin{center}
[Hội thảo Bourbaki, tập 13, 1960/61, số 221, 28 tr.]
\end{center}

\noindent\textbf{Trang 6, Định lý 2.1, dòng thứ 3.} --- Thay vì
«cần», đọc là «cần và đủ».

\noindent\textbf{Trang 15, dòng 13.} --- Thay vì «$P\neq Q$», đọc là
«$P<Q$».

\noindent\textbf{Dòng 15.} --- Thay vì «Nếu trái lại $P=Q$», đọc là
«Trong trường hợp còn lại, ta có $P(n)\geq Q(n)$ khi $n$ đủ lớn».

\noindent\textbf{Dòng 18, công thức (*).} --- Thay vì «$<$» và «$>$»,
đọc là «$\leq$» và «$\geq$».

\noindent\textbf{Dòng thứ 5 tính từ dưới lên.} --- Thay vì «$r$», đọc là
«$Z$».

\noindent\textbf{Trang 18, Nhận xét 3.9.} --- Xin lưu ý rằng Hartshorne đã
nghiên cứu các thành phần liên thông của những lược đồ Hilbert trên một trường
đóng đại số. Ông chứng minh rằng các $\operatorname{Hilb}^{P}_{r}$ là liên
thông và xác định những cặp $(r,P)$ sao cho
$\operatorname{Hilb}^{P}_{r}\neq\varnothing$\footnote{HARTSHORNE
(R.). --- Connectedness of the Hilbert scheme (Luận án, Harvard, 1962).}.

\noindent\textbf{Trang 23, dòng thứ 5 tính từ dưới lên.} --- Thay vì «một
thành phần», đọc là «của một thành phần».

\noindent\textbf{Trang 23 và 24, Nhận xét 5.5.} --- Liên quan đến ví dụ của
Zappa, xin lưu ý rằng Mumford thậm chí vừa xây dựng một thành phần bất khả quy
của lược đồ Hilbert của $\mathbf P^3_{\mathbf C}$ (các điểm tổng quát của nó
biểu diễn những đường cong không kỳ dị bậc 14, giống 24); thành phần này không
rút gọn tại điểm tổng quát của nó. Bằng cách thổi lên dọc theo một trong các
đường cong thu được, ông cũng thu được một lược đồ xạ ảnh chính quy chiều 3
trên $\mathbf C$, mà lược đồ moduli hình thức của nó không rút gọn tại điểm
tổng quát; hay tương đương, đa tạp moduli địa phương của nó, theo nghĩa của
hình học giải tích trên $\mathbf C$ (xem Hội thảo Cartan, tập 13, 1960/61),
không rút gọn tại mọi điểm.

\par
